% !TeX root = ../papers/nc-p3.tex
\section{一隅の二行二列恒等式による有限Clifford枝の内在的分類}
\label{sec:one-corner-clifford-classification}
\begin{proposition}\label{prop:all-window-corner-center}
任意の標準正則窓の標準折り畳み $B$、一つの頂点 $e$ と $A=eBe$ について、左乗法と制限は
\[
        B\simeq\operatorname{End}_{A}(Be),\qquad Z(B)\simeq Z(A)
\]
を与える。右加群の自己準同型を取り、同型は内部次数を保つ。一隅がPIなら $B$ もPIである。
\end{proposition}
\begin{proof}
$I=BeB$ と置く。$B/I$ は巡回箙から一頂点を除いた商なので有限次元である。
乗法 $Be\otimes_A eB\to I$ の核 $K$ は $Ke=0$ を満たす。実際右から $e$ を掛けると
$Be\otimes_A A\to Be$ は同型である。従って $KI=0$ で $K$ の各元が生成する右加群は有限次元である。
$B$ にはこの商の単純をsocleに持つ右部分加群はないから $\operatorname{Hom}_B(K,B)=0$。
任意の $f\in\operatorname{End}_A(Be)$ について $u\otimes v\mapsto f(u)v$ はそこで零となり、
$I\to B$ に降りる。$\operatorname{Ext}^j_B(B/I,B)=0$ ($j=0,1$) により $B\to B$ へ一意に延長し、
左乗法で $f$ を復元する。消滅は有限次元次数付き $B/I$ の頂点組成列とAS条件から従う。
非斉次準同型にも有限次元頂点加群の有限生成射影分解で同じ消滅が成立する。
$c\in Z(A)$ による $Be$ の右乗法は全ての左乗法と可換なので $Z(B)$ に復元される。
$Be$ は右 $A$ 上で有限生成である。初めて $e$ に到達する長さ3以下の経路を取ればよい。
有限生成自由右 $A$ 加群からの全射で、核を保つ自由加群自己射の部分環を取ると
$\operatorname{End}_A(Be)$ はその商である。$A$ がアフィンPIなら行列環もPIなので（行列環のPI保存）結論を得る。
\end{proof}

\begin{theorem}\label{thm:one-corner-clifford-classification}
標準有限Frobenius候補の四辺が有限とする。次は同値である。
\begin{enumerate}
\item 窓が標準Koszul左右AS正則で、一頂点の隅が二行二列行列の全多項式恒等式を満たす。
\item 窓が定理\ref{thm:index-two-central-clifford}の、基点のない四二次式系で階数2以下との交叉が
有限なClifford標準形に添字保持同型である。
\end{enumerate}
従ってこの枝には閉包、Noether性、半素性、四次中心パラメータの仮定は要らない。
任意の特殊値を含む同型は二次式四空間の $\mathrm{PGL}_4$ 軌道の一致と同値で、最小周期は全て1である。
\end{theorem}
\begin{proof}[単純の分離と中心有限性]
第一条件を仮定する。前命題から全環はPIである。非頂点単純 $S$ は全頂点に非零であり、
どの隅にもその頂点空間上で単純に作用させる。一つの頂点次元が1なら $S$ は全頂点一次元である。
実際その位置への双対の像が直線となり、有限グラフに沿う直線の伝播と四段後の固有直線の議論
（定理\ref{thm:finite-two-dimensional-quadrics}）を任意の次元で適用できる。
この薄い単純と各頂点単純との第一Extは両側とも零である。例えば一頂点を商に加える拡張の
余分な矢はその頂点を出る四つのスカラーで、次辺の点評価と関係を満たす。その評価の階数は3で、
許される四成分は元の矢の倍だけである。これは持ち上げる頂点ベクトルの変更に一致する。逆向きも前辺で同じである。

非薄単純では指定頂点の次元が少なくとも2で、行列恒等式の上界からちょうど2である。
その位置の正長ループの線型結合の像は $M_2(k)$ 全体である。正次数の隅イデアルの像が零なら
その単純次元は1であり、非零なら単純行列環のイデアルとして全体である。
二つの正次数元 $X,Y$ を $E_{12},E_{21}$ に作用するよう取ると $[X,Y]^2$ は単位に作用し、
二行二列恒等式によって隅の中心元である。その正次数の斉次部分の一つが非零作用する。
前命題でこれを全中心に延ばすと頂点と $S$ は中心指標で分離する。
従って定理\ref{thm:pi-vertex-extension-criterion}が適用でき、$B$ は正規な有限生成中心上有限で
合成整域である。一般隅の次数 $s$ は高々2で、全頂点隅で等しい。全非頂点指標上の単純の各頂点次元も
一定である。中心元が全頂点単純とのExtを消すため、四つのKoszul Euler数が全て零となり、
巡回差分の四乗の核は定数列だけだからである。
\end{proof}

\begin{proof}[四次被約traceの基点]
非頂点の閉中心指標 $\chi$ を固定する。隅の次数 $s\le2$ の被約traceは、その指標における
全単純型の行列traceの正整数重み付き和であり、次元の重み付き和は $s$ である。
定理\ref{thm:pi-vertex-extension-criterion}の安定格子と冪等元のCayley--Hamiltonの議論を
正規中心上有限なこの隅に適用すればよい。これらの隅単純は $B$ の同じ指標の単純の制限で全てを与える。
非頂点指標では $B/BeB$ の単純が存在せず、繊維の $e$ は充満冪等元となるからである。

二次元単純があればそれだけが重み1で現れる。その隅の長さ4の語の像の線型包は $M_2$ である。
実際各矢行列空間 $H,H'$ は次元3以上で、その直交には階数1の行列がない。
隣接積の零化に $D\ne0$ があれば $DH'\subset H^\perp$。$D$ が可逆なら次元で不可能、
階数1ならその全像は階数1以下で $DH'=0$ が必要になるが $H'$ は可逆行列を含む。
よって二段で既に全行列環を張り、後の二段にも可逆な矢の線型結合がある。
従ってtraceが非零の長さ4の経路線型結合がある。

薄い単純だけのとき、長さ4のtraceはその各点の四文字の評価テンソルの重み付き和である。
異なるグラフの台点なら二つの純テンソルは比例しないので消えない。一つの単純型だけなら
標数零で正の重みを掛けた非零評価である。残るのは同じグラフ点 $p$ に対する異なる一周倍率の
二単純であるが、これだけの二型が一つの中心繊維をなすことはない。

実際同じ $p$ の矢 $p_i$ と $\lambda_i p_i$ の間の拡張の非対角矢を $c_i\in V_i^*$ と置く。
関係の非対角成分は $(\lambda_{i+1}c_i,c_{i+1})$ が点対応の双線型方程式の一次式を消すこととなる
（順番は一方の表現を部分とする向きで選ぶ）。$p_i$ の倍を除いた接成分は、従って
$\bar c_{i+1}=\lambda_{i+1}d\sigma_i(\bar c_i)$ を満たす。ここで $\sigma_i$ は局所点スキームの後続同型で、
定理\ref{thm:finite-cubic-separators}により四回の合成は非被約部分を含め恒等である。
$\prod_i\lambda_i\ne1$ なので全接成分が零になる。残る各 $p_i$ の倍率は頂点の四基底の
非対角変更で全て取り除ける。その巡回一次系の行列式は符号を除き $1-\prod\lambda_i$ だからである。
逆順でも同じで両向き第一Extが消える。よってこの二単純だけを持つ有限次元中心繊維は二ブロックに分かれる。
しかし閉点完備化に二つの冪等元を持ち上げると、それらの間の両方の隅は剰余零なのでNakayamaにより零である。
すると完備した環の中心に非自明冪等元があることになる。有限な中心加群の交換子の核は平坦な完備化と可換で、
完備中心は局所環であるから不可能である。

以上で全ての非頂点指標に非零となる経路長4の中心traceがある。射影中心上で四次中心系は基点がなく、
一般四元で切れば剰余は有限次元である。定理\ref{thm:successively-central-cutoff}から四次中心切断を得る。
定理\ref{thm:central-cut-center-dichotomy}で $s=2$ となり、定理\ref{thm:index-two-central-clifford}が第二条件、
例外を含む同型、周期を与える。逆にその形では既に正則で一般隅が次数2の可除環なので指定した恒等式を満たす。
\end{proof}

\begin{corollary}\label{cor:no-index-three-finite-noetherian-window}
左右Noetherな正則有限四辺では、一隅が三行三列行列の多項式恒等式を満たす場合も上と同じClifford枝に尽きる。
特に中心有限な正則有限四辺に頂点可除指数3はない。
\end{corollary}
\begin{proof}
一隅から全環がPIとなり定理\ref{thm:noetherian-pi-window-cutoff}で中心上有限、一般隅の次数は $s\le3$ である。
四次traceの基点の証明に追加される分割は $3$、$2+1$、および薄いもの三つだけである。
各頂点三次元の単純では各非零ベクトルで矢像は二次元以上、双対も同じである。従ってどの二平面の像も全空間を張り、
隣接二辺の積の張空間は各非零ベクトル上で全射となる。そのtrace直交は階数1行列を含まず、
三行三列のSegreの余次元4より直交の次元は4以下、積空間の次元は5以上である。
反対の積空間も次元5以上なので九次元内で互いに直交できず、四次traceは非零である。
$2+1$ の場合、二次元単純のtraceテンソルの隣接二対での行列階数は4である。両積空間がともに $M_2$ で
trace対合が完全であることによる。薄いものの階数1との和は零にならない。
全てが薄い場合、今現れる高々三つの異なる台点の四文字評価は独立である。これらに支えられた関係は長さ20のグラフ上の退化した積対合を与え、定理\ref{thm:finite-four-slot-rank}で排除される。
和が零なら各台点には少なくとも二つの異なる倍率型が必要で、重み総和3以下では一台点しかない。
しかし異なる倍率型の間は全て上の拡張消滅で分離し、一つの中心繊維のブロック連結性に反する。
よってこの場合も四次四元が取れ、中心切断の二択と $s\le3$ から $s=2$ となる。
\end{proof}
